%%%%%%%%%%%%%%%%%%%%%%%%%%%%%%%%%%%%%%%%%%%%%%%%%%%%%%%%%%%%%%%%%%%%%%%%%%%%%%%%
%%
\cleardoubleoddpage%  Make sure to start each chapter on a new odd page
\chapter{Discussion}
\label{ch:discussion}


\section{Interpretation of Results}
\label{sec:discussion:interpretation}

% TODO: broaden back out from the technical results. Which components
% contributed most? Did the moving mesh help on slowly evolving pathology,
% or was uniform mesh enough? Did the surrogate encoder improve
% generalization across patients?


\section{Clinical Implications}
\label{sec:discussion:clinical}

% TODO: which prediction horizon is reliable? What RMSE/Dice level would be
% clinically useful? Where does the model fail, and would those failures
% matter in practice?


\section{Limitations}
\label{sec:discussion:limitations}

% TODO: data quantity (sample size, missing visits); normalization of the
% binary mask channel; monitor function chosen by heuristic rather than
% derived for vector-valued states; fixed-grid assumption discards
% peripheral retina; encoder sees predicted (not observed) layers during
% pushforward unrolling.


\section{Future Work}
\label{sec:discussion:future}

% TODO: incorporating surrogate models of the next state into DMM's monitor
% function (the future-state problem from MM-PDE's conclusion); extending
% to 3D volumetric OCT; multi-task joint training across multiple retinal
% pathologies; uncertainty quantification for clinical deployment.


%%
%%%%%%%%%%%%%%%%%%%%%%%%%%%%%%%%%%%%%%%%%%%%%%%%%%%%%%%%%%%%%%%%%%%%%%%%%%%%%%%%
